\begin{abstract}
Scientific workflow schedulers reserve memory before execution, balancing
insufficient allocations against wasted capacity. Consumption-Aware Memory
Prediction (\toolname), introduced in our WORKS paper, uses completed tasks'
consumption and memory observations to inform later requests. We examine
its deployment through measurement and feature ablations, downstream
selective-auditing analysis, auditing-cost estimation, and live Bowtie2
experiments. The live study evaluates repeated inputs and transfer to a
separate input study, with 6,000 tasks per stage. Repeating the original
inputs reduces aggregate requests by 75.34\%, while unused
memory-time falls from 32.86089 to 0.01847~GiB-hours. On new inputs, continuing
from this warm state requests 319.22~GiB with 0.12805~GiB-hours of
unused memory-time; continuing directly from the shared cold
checkpoint requests 343.56~GiB with 0.22109~GiB-hours of unused memory-time.
Both branches keep learning online.
Together, the experiments characterize the benefits, measurement costs,
and limits of consumption feedback under repeated and previously unseen
workflow inputs.
\end{abstract}

\begin{IEEEkeywords}
scientific workflows, online memory allocation, consumption feedback,
eBPF, selective auditing
\end{IEEEkeywords}
